Several approaches have been proposed for combining reinforcement learning (RL) with spiking neural networks (SNNs). Early work focused on biologically plausible learning rules, such as Spike-Timing-Dependent Plasticity (STDP), which have been successfully applied to tasks like source seeking \cite{Florian2005} and maze navigation \cite{fremaux2013spiking}. While these methods demonstrate the potential of SNNs for solving control tasks, they often lack the scalability and efficiency of modern deep RL techniques. 

To benefit from the advances in deep RL, surrogate gradients can be used to enable end-to-end training of SNNs with standard RL algorithms. The Deep Spiking Q-Network (DSQN) algorithm \cite{Chen2022} adapts the DQN algorithm \cite{mnih2013playingatarideepreinforcement} to the spiking domain. Although DSQN achieves robust performance and demonstrates significant energy savings on neuromorphic hardware \cite{akl2021porting}, it is trained on single transitions and resets the network state at every step Therefore, it resets the network state at every environment transition, preventing the model from capturing temporal dependencies across time. Moreover, as a value-based method, DSQN is inherently limited to discrete action spaces and is not well-suited for continuous control.

More generally, value-based RL algorithms that leverage recurrent architectures, such as R2D2~\cite{Kapturowski2019} rely on a warm-up period to stabilize hidden states before applying updates. This technique assumes the agent can gather long sequences of experience. However, in robotic control environments, such as drone flight, poorly initialized policies often result in early termination (e.g., due to crashing), preventing the network from bridging the warm-up period. This limits the practicality of recurrent value-based methods in high-risk real-wrold environments.

To address these limitations, policy-based approaches have been explored. SNN-PPO \cite{zanatta2024exploring} applies the Proximal Policy Optimization algorithm to train SNNs in an on-policy manner using entire trajectory sequences. This method has been shown to successfully solve a variety of continuous control tasks from the MuJoCo suite \cite{mujoco}. On-policy methods like SNN-PPO can be challenging to use when sample cost is high. 

Off-policy approaches enable reusing past experience stored in the replay buffer. PopSAN \cite{Tang2020} is an off-policy, actor-critic method designed for continuous control with SNNs. By leveraging a conventional ANN as the critic, PopSAN benefits from the stability of standard deep RL while using an SNN-based actor. It achieves strong performance on several MuJoCo tasks in simulation and reports up to a 140x reduction in energy consumption compared to non-neuromorphic implementations \cite{Tang2020_Hardware}. Nonetheless, similar to DSQN, PopSAN resets the hidden states between transitions, limiting its ability to capture temporal dependencies across timesteps.

